% What a closer look is and how it is made: the same for every model. A paper
% reads it under its own section heading, before the cuts that are its own.
\iflook
A leaderboard says which run is best. It does not say where a model is right
and where it is wrong, and one score is a poor guide to whether the app can
use it. So the same test rows are cut several ways, on the machine that has
the dataset and every run's test predictions, and the small tables that come
out are committed beside the leaderboard. Everything in this section and the
next is drawn from them (\V{look.test-rows} test rows, looked at on
\V{look.date}).

\paragraph{Slices.}
The test rows are grouped by one column at a time. For each group a slice
gives how many rows fall in it, what happened there, what a run said on
average, and how the run scored. Two runs are set side by side: \V{challenger.name},
the best run that is not the reference, and the reference. The gain of the
first over the second is measured inside each group as it is for the whole
test set, paired and over resampled games (\V{look.resamples} draws), and is
left out for a group with fewer than \V{look.smallest-compared} test rows.
A slice is where a claim such as ``the network helps most on long lines'' is
checked, not made: a group whose interval holds zero says nothing either way.

\paragraph{Grids.}
Two columns at once: what happened in each cell, and what the run said there,
on one scale of colour so the two can be compared by eye. A cell with fewer
than \V{grid.smallest-cell} test rows is left empty.

\paragraph{Examples.}
Positions are chosen by rule, never by eye: ordinary rows picked by a hash of
their key, the rows the best run is surest of at each end, and its largest
misses each way, \V{examples.count} of each. They show what the numbers are
about. They are not evidence of how often anything happens.
\else
\emph{This build of the paper was drawn before the model's closer look was
run. The closer look cuts the test rows by rating, clock and the shape of a
position, checks what a stated number is worth, and shows example positions.
It is written by \texttt{lab analyse} on the machine that holds the dataset
and every run's test predictions, and this section and the next appear once
its file is committed.}
\fi
